% Part: sets-functions-relations
% Chapter: relations
% Section: reflections
%
\documentclass[../../../include/open-logic-section]{subfiles}

\begin{document}

\olfileid[tr]{sfr}{rel}{ref}
\olsection{Felsefi Düşünceler}

\olref[set]{sec} kesiminde bağıntıları belirli kümeler olarak tanımladık. Burada
durup kısa bir felsefi soru sormalıyız: böyle bir tanım ne \emph{yapmaktadır}?
Birtakım metafizik özdeşlik olgularını \emph{keşfettiğimizi} söylemek
isteyeceğimizden ciddi biçimde kuşku duymalıyız; örneğin $\Nat$ üzerindeki sıralama
bağıntısının, \olref[set]{sec} kesiminde tanımladığımız
$R=\Setabs{\tuple{n,m}}{n, m \in \Nat\text{ ve } n < m}$ kümesi olduğunun
\emph{ortaya çıktığını} söylemek istemeyiz. Bunun üç nedeni vardır.

Birincisi: \olref[set][pai]{wienerkuratowski} içinde sıralı ikiliyi
$\tuple{a, b}=\{\{a\}, \{a, b\}\}$ olarak tanımladık. Bunun yerine
$\lVert a,b\rVert = \{\{b\}, \{a, b\}\}=\tuple{b,a}$ tanımını düşünün.
$a \neq b$ olduğunda $\tuple{a, b} \neq \lVert a,b\rVert$ olur. Ancak sıralı
ikili tanımımız olarak $\tuple{a,b}$ yerine $\lVert a,b\rVert$ ifadesini de
aynı ölçüde benimseyebilirdik. Her iki tanım da eşit ölçüde işe yarardı. Böylece
doğal sayılar üzerindeki sıralama bağıntısı ``olmaya'' eşit ölçüde uygun iki
adayımız vardır:
\begin{align*}
		R &= \Setabs{\tuple{n,m}}{n, m \in \Nat \text{ ve }n < m}\\
		S &= \Setabs{\lVert n,m\rVert}{n, m \in \Nat \text{ ve }n < m}.
\end{align*}
$R \neq S$ olduğuna göre, kaplamsallık uyarınca ikisinin birden $\Nat$ üzerindeki
sıralama bağıntısıyla özdeş olamayacağı açıktır. Öte yandan olgusal bir mesele
olarak $S$ yerine $R$'nin (ya da $R$ yerine $S$'nin) sıralama bağıntısının kendisi olduğunu
öne sürmek keyfî ve dolayısıyla biraz mahcup edici olurdu. (Bu, Benacerraf'ın
\citeyear{Benacerraf1965} tarihli çalışmasında ün kazandırdığı küme kuramsal
indirgemeciliğe karşı savın çok basit bir örneğidir. Bu konuya birkaç kez
yeniden döneceğiz.)

İkincisi: \emph{her} bağıntının bir kümeyle özdeşleştirilmesi gerektiğini
düşünüyorsak, kümeye eleman olma bağıntısı $\in$ de belirli bir küme olmalıdır.
Gerçekten de bunun $\Setabs{\tuple{x,y}}{x \in y}$ kümesi olması gerekirdi.
Peki bu küme var mıdır? Russell paradoksu göz önüne alındığında böyle bir
kümenin var olduğu hiç de önemsiz bir iddia değildir. Nitekim
\oliflabeldef{cumul:::part}{\olref[cumul][][]{part} kısmında geliştirdiğimiz
küme teorisi bu kümenin varlığını \emph{reddedecektir}.\footnote{Biraz ilerisini
söylersek, gerekçe şudur: olmayana ergi yoluyla
$I=\Setabs{\tuple{x,y}}{x \in y}$ kümesinin var olduğunu varsayalım. O zaman
$\bigcup \bigcup I$ evrensel kümedir; bu da
\olref[sfr][z][sep]{thm:NoUniversalSet} ile çelişir.}}{küme teorisini aksiyomatik
bir teori olarak titizlikle geliştirmek mümkündür ve bu teori gerçekten de bu
kümenin varlığını reddeder.} Dolayısıyla bazı bağıntılar küme olarak ele
alınabilse bile, kümeye eleman olma bağıntısı özel bir durum olmak zorundadır.

Üçüncüsü: bağıntıları kümelerle ``özdeşleştirirken'', $\tuple{x,y} \in R$
yerine $Rxy$ yazmaya izin vereceğimizi söylemiştik. Elemanlık bağıntısı
``$\in$'' bir yüklem \emph{olarak} ele alındığı sürece bunda sorun yoktur. Ancak
``$\in$'' işaretinin belirli türde bir kümeyi gösterdiğini düşünürsek,
``$\tuple{x,y} \in R$'' ifadesi yalnızca kümeleri gösteren üç tekil terimden
oluşur: ``$\tuple{x,y}$'', ``$\in$'' ve ``$R$''. Böyle bir ad listesi,
``fincan kalemlik masa'' anlamsız dizisinin bir önerme dile getirebilmesinden
daha fazla önerme dile getiremez. Yine, bazı bağıntılar kümeler olarak ele
alınabilse bile elemanlık bağıntısı özel bir durum olmak zorundadır. (Burada
Frege'nin \emph{kavram at} paradoksunun basit bir biçimiyle Wittgenstein'ın bir
zamanlar Russell'a yönelttiği ünlü itiraz bir araya getirilmektedir.)

Öyleyse sonuç nedir? Sayılar üzerindeki bağıntıların kümeler olduğunu
söylememizde \emph{yanlış} bir şey yoktur. Yalnızca bu sözün hangi anlayışla
söylendiğini kavramamız gerekir. Metafizik bir özdeşlik olgusunu dile
getirmiyoruz. Yalnızca belirli bağlamlarda (belirli) bağıntıları belirli kümeler
\emph{olarak ele alabileceğimizi} (ve alacağımızı) belirtiyoruz.

\end{document}
